% Add "answers" option to print the answers
\documentclass[addpoints]{exam}

\usepackage[utf8]{inputenc}
\usepackage{calc}
\usepackage[dvipsnames]{xcolor}
\usepackage{ifthen}
\usepackage{tcolorbox}

% \makebox[\textwidth]{Name:\enspace\hrulefill}
%  
% \vspace{5mm}
%\makebox[\textwidth]{Instructor’s name:\enspace\hrulefill}

\definecolor{MyLightGray}{gray}{0.60}

\newcommand{\truefalse}
{		\begin{checkboxes}
			 \choice True
			 \choice False. Explanation: 
		\end{checkboxes}
		\vspace{10pt}
}


% Check if 'answers' is an option
\newcommand{\ifanswers}[2]{%
  \ifthenelse{\boolean{exam@answers}}{#1}{#2}%
}



\newcommand{\answerspace}[1]
{
	\ifprintanswers
	\else
		\setlength{\mylength}{10pt}%
		% \vspace{#1\mylength}
		% Make the grey box
		\par\nobreak\noindent
		\makebox[\linewidth][l]{%
			% set inner padding to zero, so only the frame shows
			{\setlength{\fboxsep}{0pt}%
			% draw a color frame, white (or transparent) inside
			\fcolorbox{MyLightGray}{white}{%
			% the inner box has zero width content but fixed height
			\parbox[t][#1\mylength][c]{%
				% subtract twice the frame‐thickness from the full width
				\dimexpr\linewidth-2\fboxrule\relax
			}{\strut}%
			}%
			}%
		}%
		\par
		\vspace{3\mylength}

  	\fi
}
\newlength{\mylength}

\newcommand{\code}[1]{{\texttt{#1}}}

\newcommand{\answer}[1]{\begin{solution}{#1}\end{solution}}




%\usepackage{listings}
\usepackage{color}
\usepackage{lastpage}
\usepackage{verbatim}
\usepackage{fancyvrb}
\usepackage{graphicx}
\usepackage{float}
\usepackage{minted}

% For those sweet, sweet acmart fonts.
\usepackage[T1]{fontenc}
\usepackage[tt=false, type1=true]{libertine}
\usepackage[varqu]{zi4}                  % Inconsolata for \ttfamily

\newcommand{\Code}[1]{\mintinline{python}{#1}}

\renewcommand{\solutiontitle}{\noindent\textbf{Solution sketch: }}

\definecolor{SolutionColor}{rgb}{0.8,0.9,1}

\cfoot{Page \thepage\ of \pageref{LastPage}}

\rhead{Initials: \rule{20mm}{0.4pt}}

\begin{document}
\fvset{numbers=left,numbersep=3pt}
\thispagestyle{empty}

\begin{center}
	{\LARGE CS 408 Fall 2026 -- Midterm 2 \\}

\fbox{\fbox{\parbox{5.5in}{\centering
	\begin{itemize}
        \item The exam is closed notes, closed book.

        \item The exam should be completed individually: any
          communication with your fellow students is \textbf{strictly}
          prohibited.

        \item If you have a question, raise your hand and the
          course staff will come to you.

        \item Please print or write legibly.

        \item Try to be concise: when a question asks for an
          explanation, a few sentences should suffice. Aim for
          quality, not quantity!

        \item Show enough of your reasoning so that we can award
          partial credit for a correct approach.

        \item Answer the questions in the spaces provided. Your answer
          does not need to take up all the space provided (hopefully
          we've given you more than enough room!). If you do run out
          of space, make a note and then use the empty pages at the
          back of the exam to complete your answer. Be sure to clearly
          label any answers on those pages with the corresponding
          question.

        \item Several questions ask you to provide something ``if
          possible.'' If it is not possible, you should plainly say so
          and then explain why that is the case.

        \item Coding questions should be answered in Python, unless
          stated otherwise.

        \item Unless otherwise noted, CFG nodes represent basic
          blocks. When necessary, refer to nodes using the labels
          given in the figure.

        \item Please double check that your copy of the exam
          is complete; you should have \pageref{LastPage}
          pages.

        \item You have 75 minutes to complete the exam.
        \end{itemize}
	}}}
\end{center}

\begin{center}
	\makebox[0.6\textwidth]{
		\combinedgradetable[h][questions]
	}
\end{center}
\vspace{20mm}

\makebox[\textwidth]{\textbf{Name (please print)}: \rule{100mm}{0.8pt} }

\vspace{10mm}
\makebox[\textwidth]{\textbf{PUID}: \rule{40mm}{0.8pt} \hspace{35mm}}

\vspace{10mm}
\makebox[\textwidth]{\textbf{Ticket \#}: \rule{40mm}{0.8pt} \hspace{35mm}}

\newpage
\begin{center}
  [PAGE DELIBERATELY LEFT BLANK]
\end{center}

\begin{questions}

  \newpage

% ========================================================
% Q1: Warm up questions (Can we pull variants of some of the simpler
% questions from in-class quizzes?)
% ========================================================

\question Short answers.
\begin{parts}

\part[2] Every property-based testing framework consists of three key
  components. List each of those pieces and explain in a few words the
  role of each.
  \begin{solutionorbox}[1.5in]
    A generator that produces inputs, a property that must hold for
    every input, and a shrinker that reduces a failing input to a
    small one.
  \end{solutionorbox}

\part[1] What is one advantage of a grammar-based fuzzer over a
  mutation-based fuzzer?
  \begin{solutionorbox}[1.5in]
    Every input generated by a grammar is valid by construction, so
    none of the budget is wasted on inputs the parser rejects.
  \end{solutionorbox}

\part[1] What is one advantage of a mutation-based fuzzer over a
  grammar-based fuzzer?
  \begin{solutionorbox}[1.5in]
    Mutation: users do not need to specify the grammar, users only
    need to supply some input seeds, it can reach malformed inputs
    that the grammar cannot express, etc.
  \end{solutionorbox}

\part[1] Suppose a symbolic execution engine hands off a path
  constraint to an SMT solver and gets \Code{UNSAT} in response (i.e.,
  the solver could not find a model). What does this tell us about the
  corresponding test path?
  \begin{solutionorbox}[1.5in]
    The corresponding path is infeasible!
  \end{solutionorbox}
\end{parts}

\newpage
\begin{center}
  [PAGE DELIBERATELY LEFT BLANK]
\end{center}
\newpage

% ========================================================
% Q2+3: Test Oracles [from spring 2026 final]
% ========================================================

\question[2] Which category of test oracle discussed in class would an
  off-the-shelf blackbox fuzzer typically use? Explain your answer in
  2-3 sentences.

  \begin{solutionorbox}[3.5in]
    An implicit oracle is the standard, because it works 'out of the
    box' on an arbitrary program, and greybox fuzzers are typically
    designed to be as push-button as possible.
  \end{solutionorbox}

\question[3] In 2-4 sentences, describe an automated testing scenario
  where a blackbox fuzzer might use one of the other kinds of test
  oracle we have discussed in class. Be sure to explicitly state which
  category of test oracle is used, and to justify why it is reasonable
  to use in this scenario.

  \begin{solutionorbox}[4in]
    There are many possible answers here. The key thing is that they
    need to identify a particular application which enable us to use a
    more intelligent test oracle: e.g. when fuzzing an image
    processing algorithm we can use a metamorphic relation as an
    oracle, or when fuzz testing a compiler optimization we can use a
    reference implementation as an oracle.
  \end{solutionorbox}

\newpage
\begin{center}
  [PAGE DELIBERATELY LEFT BLANK]
\end{center}
\newpage

% ========================================================
% Q4: Delta Debugging [from spring 2026 final]
% ========================================================

\question Suppose that we have the test input
  $[1, 2, 3, 5, 3, 8, 7, 10, 42, 11]$, and a sorting program that
  fails if and only if it is given a list with at least two elements,
  one of which is the number $3$, and which contains at least one pair
  of adjacent elements where the first element is greater than the
  second. Identify which of the following test inputs could be
  constructed by the delta debugging algorithm (\Code{ddmin}) while
  minimizing this input, and whether each could be its final
  answer. (Indicate your answer by circling your choice.)

  \begin{parts}
  \part[2] $[1, 2, 3, 5, 3]$
    \begin{itemize}
    \item[ ] Constructed by \Code{ddmin}: Yes \quad/\quad No %
      \hfill Final output of \Code{ddmin}: Yes \quad/\quad No
    \end{itemize}

  \part[2] $[1, 3, 8, 10]$
    \begin{itemize}
    \item[ ] Constructed by \Code{ddmin}: Yes \quad/\quad No %
    \hfill Final output of \Code{ddmin}: Yes \quad/\quad No
  \end{itemize}

  \part[2] $[3, 7, 8]$
    \begin{itemize}
    \item[ ] Constructed by \Code{ddmin}: Yes \quad/\quad No %
    \hfill Final output of \Code{ddmin}: Yes \quad/\quad No
  \end{itemize}

  \part[2] $[5, 3]$
    \begin{itemize}
    \item[ ] Constructed by \Code{ddmin}: Yes \quad/\quad No %
    \hfill Final output of \Code{ddmin}: Yes \quad/\quad No
  \end{itemize}

\part[2] $[3, 6]$
    \begin{itemize}
    \item[ ] Constructed by \Code{ddmin}: Yes \quad/\quad No %
      \hfill Final output of \Code{ddmin}: Yes \quad/\quad No
    \end{itemize}
  \end{parts}

\question Recall that the coverage-guided fuzzer from the project
  minimized certain inputs, but did not use \Code{ddmin} to do so.

  \begin{parts}
  \part[2] In a couple of sentences, compare and contrast one relative
    advantage and one relative disadvantage of \Code{ddmin} versus one
    of the other minimization strategies we have seen in class. (You
    do not necessarily have to discuss the shrinking strategy from the
    project.)
    \begin{solutionorbox}[2in]
      The project used a byte-by-byte shrinking strategy guaranteed to
      be linear in the size of the input, which is faster than
      \Code{ddmin} in the worst case. The downside is that the project's
      strategy is not guaranteed to produce a 1-minimal example.
    \end{solutionorbox}

  \part[2] In 1-2 sentences, what did the project's shrinker check for
    when deciding whether to trim a particular byte from an input and
    why?
    \begin{solutionorbox}[1.5in]
    It checked whether the trimmed input still produced the same coverage.
  \end{solutionorbox}
\end{parts}

\newpage
\begin{center}
  [PAGE DELIBERATELY LEFT BLANK]
\end{center}
\newpage

% ========================================================
% Q3: Coverage guided fuzzing (maybe we can say: given these 3-4
% inputs, identify each as being produced by a grey box fuzzer using
% these mutators and seed set, and whether or not it would be included to the seed
% input after running.
% ========================================================

\question Suppose that a grey-box fuzzer is testing the following
  function using the input \Code{"good"} as its sole initial seed
  value:

\begin{minted}{python}
def crashme(s: str) -> None:
    if len(s) > 0 and s[0] == 'b':
        if len(s) > 1 and s[1] == 'a':
            if len(s) > 2 and s[2] == 'd':
                if len(s) > 3 and s[3] == '!':
                    raise Exception()
\end{minted}

  The fuzzer is equipped with the following mutation operations:
  \begin{itemize}
  \item delete a character at random,
  \item replace a character with another randomly chosen character, or
  \item append a space to the end.
  \end{itemize}

  At each iteration, the fuzzer is allowed to apply either \textbf{one
    or two} of these mutations. For each of the following inputs,
  indicate whether the fuzzer could \textbf{produce} it from its
  initial seed; if so, you should additionally indicate whether the
  fuzzer would then \textbf{include} it in the population. (Indicate
  your answer by circling your choice.)

\begin{parts}

\part[2] \Code{"bood"}
  \begin{itemize}
  \item[ ] Produced from \Code{"good"}: Yes \quad/\quad No %
    \hfill Added to the population: Yes \quad/\quad No
  \end{itemize}

\part[2] \Code{"cood"}
  \begin{itemize}
  \item[ ] Produced from \Code{"good"}: Yes \quad/\quad No %
    \hfill Added to the population: Yes \quad/\quad No
  \end{itemize}

\part[2] \Code{" good"}
  \begin{itemize}
  \item[ ] Produced from \Code{"good"}: Yes \quad/\quad No %
    \hfill Added to the population: Yes \quad/\quad No
  \end{itemize}

\part[2] \Code{"baod"}
  \begin{itemize}
  \item[ ] Produced from \Code{"good"}: Yes \quad/\quad No %
    \hfill Added to the population: Yes \quad/\quad No
  \end{itemize}

\part[2] \Code{"bad!"}
  \begin{itemize}
  \item[ ] Produced from \Code{"good"}: Yes \quad/\quad No %
    \hfill Added to the population: Yes \quad/\quad No
  \end{itemize}

  % \begin{solutionorbox}[2.6in]
  %   The seed \Code{"good"} reaches only the first \Code{if}. An input
  %   is added when it covers something new.

  %   \begin{center}\small
  %     \begin{tabular}{llll}
  %       & produced? & added? & \\ \hline
  %       \Code{"bood"}  & yes, 1 edit  & \textbf{yes} & change \Code{g} to \Code{b}; reaches the second \Code{if} \\
  %       \Code{"cood"}  & yes, 1 edit  & no           & change \Code{g} to \Code{c}; covers nothing new \\
  %       \Code{" good"} & \textbf{no}  & ---          & a space may only be appended, not prepended \\
  %       \Code{"baod"}  & yes, 2 edits & \textbf{yes} & \Code{g}$\to$\Code{b}, \Code{o}$\to$\Code{a}; reaches the third \Code{if} \\
  %       \Code{"bad!"}  & \textbf{no}  & ---          & four edits away from \Code{"good"} \\
  %     \end{tabular}
  %   \end{center}
  % \end{solutionorbox}

\part[2] For the inputs which could not be produced from the initial
  seed, could the fuzzer eventually consider them? Explain why or why
  not in 2-3 sentences.
  \begin{solutionorbox}[2in]
    Each input that covers something new is added to the population
    and becomes a new starting point for mutation. So the fuzzer does
    not have to apply four edits at once: it reaches \Code{"bood"},
    keeps it, mutates \emph{that} to reach \Code{"baod"}, keeps it,
    and so on. The coverage feedback is what tells it which
    intermediate inputs are worth keeping --- a black-box fuzzer
    discards all of them and never gets past the first character.
  \end{solutionorbox}
\end{parts}

\newpage
\begin{center}
  [PAGE DELIBERATELY LEFT BLANK]
\end{center}
\newpage

% ========================================================
% Q4: Symbolic execution
% ========================================================

\question Consider the following function:

\begin{minted}{python}
def check(x: int, y: int) -> int:
    r: int = 0
    s: int = 0
    if x > 0:
        r = 1
    if y > 0:
        if x < 0:
            r = r + 10
        else:
            r = r + 20
    return (r, s)
\end{minted}

\begin{parts}

\part[6] Draw the symbolic execution tree that a symbolic execution
  engine would produce for \Code{check}. Use $a$ and $b$ for the
  initial symbolic value for \Code{x} and \Code{y}, respectively, so
  that the initial symbolic state is:
  $\{\texttt{x}\mapsto a, \texttt{y} \mapsto b\}$. When drawing the
  tree, you do not need to explicitly annotate nodes with symbolic
  states.

  \begin{solutionorbox}[4in]
    Branching on \Code{x > 0}, then on \Code{y > 0} beneath each, and
    then on \Code{x < 0} beneath each \Code{y > 0}. The
    \Code{y <= 0} branches are leaves.

    \begin{center}\small
      \begin{tabular}{cl}
        leaf & path condition \\ \hline
        1 & $a > 0 \wedge b > 0 \wedge a < 0$ \\
        2 & $a > 0 \wedge b > 0 \wedge \neg(a < 0)$ \\
        3 & $a > 0 \wedge \neg(b > 0)$ \\
        4 & $\neg(a > 0) \wedge b > 0 \wedge a < 0$ \\
        5 & $\neg(a > 0) \wedge b > 0 \wedge \neg(a < 0)$ \\
        6 & $\neg(a > 0) \wedge \neg(b > 0)$ \\
      \end{tabular}
    \end{center}

    Note that the third \Code{if} is nested inside the second, so it is
    only reached on the \Code{b > 0} side. That is why there are six
    leaves and not $2^3 = 8$.
  \end{solutionorbox}

\part[1] How many test paths does this tree represent?
  \begin{solutionorbox}[1in]
    The tree has \textbf{six} leaves, so it represents six test paths.

    Only \textbf{five} of them can be executed. Leaf 1 requires
    $a > 0 \wedge a < 0$, which no value of \Code{a} satisfies, so that
    path is \textbf{infeasible}.

    (Leaf 5 is feasible, but only just: $\neg(a > 0) \wedge \neg(a < 0)$
    forces $a = 0$.)
  \end{solutionorbox}

\newpage
\begin{center}
  [PAGE DELIBERATELY LEFT BLANK]
\end{center}
\newpage

\part[2] Now suppose we extend \Code{check} to return a pair of
  values, adding two more conditionals at the end of the function to
  compute the second component:

\begin{minted}{python}
def check(x: int, y: int, z: int) -> tuple:
    r: int = 0
    s: int = 0
    if x > 0:
        r = 1
    if y > 0:
        if x < 0:
            r = r + 10
        else:
            r = r + 20
    if z > 0:
        s = 1
    if z % 2 == 0:
        s = s + 2
    return (r, s)
\end{minted}

  How many leaves would the symbolic execution tree have now?
  (\textbf{Note:} you do not need to draw the new tree!) Which (if
  any) of the practical limitations of symbolic execution that we
  discussed in class does this point to?
  \begin{solutionorbox}[3in]
    The two new conditions are independent of each other and of
    everything above them, so every existing path splits four ways:
    \[
      6 \times 2^2 = \textbf{24 leaves.}
    \]

    Of these, \textbf{20} are feasible. The single infeasible prefix
    from part (b) now accounts for \emph{four} unreachable leaves
    rather than one --- which is why an engine checks feasibility as it
    descends rather than waiting until it reaches a leaf.

    More generally, each additional independent branch \textbf{doubles}
    the number of paths. A realistic program has thousands of branches,
    so the number of paths is astronomically large and no engine can
    enumerate them all. This is the \textbf{path explosion} problem,
    and it is why practical tools bound the search instead.
  \end{solutionorbox}

\end{parts}

\newpage
\begin{center}
  [PAGE DELIBERATELY LEFT BLANK]
\end{center}
\newpage

% ========================================================
% Q5: Concolic execution [from spring 2026 final]
% ========================================================

\question Consider the following program: \vspace{-.3cm}
  \begin{minted}{python}
def bar(arr: list) -> int:
    s: int = 0
    if len(arr) < 3:
        return -1
    for i in range(3):
        if arr[i] >= i:
            s = s + arr[i]
        else:
            s = s - arr[i]
    if s % 2 == 0:
        return 1
    else:
        return 0
\end{minted}
\begin{parts}

\part[6]
  List the final path constraint that concolic execution would produce
  for the program \Code{bar} on the test input %
  \Code{bar([1, 0, 3])}. Give your answer as a sequence of path
  conditions separated by $\land$, e.g., for the path conditions $P$,
  $Q$, and $R$, write $P \land Q \land R$ . When writing your path
  condition, assume the initial symbolic input has the form
  $arr \mapsto [a , b, c]$, and write the symbolic value stored at
  index \Code{0} as $a$, \Code{1} as $b$, and \Code{2} as $c$. You may
  drop clauses that are always true, e.g., $3 < 5 + 1$ from the path
  constraint, but otherwise you should not simplify it in any way.
  \begin{solutionorbox}[2in]
    The final path constraint is:\\
    $\lnot$ \Code{len(arr) < 3} $\land$ \Code{a >= 0} $\land$ $\lnot$
    \Code{b >= 1} $\land$ \Code{c >= 2} $\land$ \Code{(a - b + c)} \%
    \Code{2 == 0}
    \end{solutionorbox}
\part[4] List the next two new inputs that a concolic execution engine
  might generate after executing \Code{bar} on the test input
  $[1, 0, 3]$. Explicitly state the strategy you used. %
  \begin{solutionorbox}[2.75in]
    Anything that flips the last path constraint they had above will
    work. Assuming they correctly generated the path condition in the
    previous part, we need an array with three elements whose first
    element is greater than -1, second element is less than 1, third
    element is greater than 1, and whose first three elements sum up
    to an odd number, e.g.,  $arr \mapsto$ \Code{[1, 0, 2]}. \\
    For the second input, we need an array whose third element is less
    than 2, e.g., $arr \mapsto$ \Code{[1, 0, 1]}.
  \end{solutionorbox}
\end{parts} %
%
\newpage
\begin{center}
  [PAGE DELIBERATELY LEFT BLANK]
\end{center}
\newpage

% \newpage
% \begin{center}
%   [PAGE DELIBERATELY LEFT BLANK]
% \end{center}
% \newpage

\if 0 %%%%%%%%%%%%%%%%%%%%%%%%%%%%%%%%%%%%%%%%%
	\question[10]
	\answerspace{30}
	\begin{solution}
        \end{solution}
\fi %%%%%%%%%%%%%%%%%%%%%%%%%%%%%%%%%%%%%%%%%

	\newpage
	\begin{center}
		[EMPTY PAGE 1 -- Use as overflow to complete your
                answers ]
	\end{center}

	\newpage
	\begin{center}
		[EMPTY PAGE 2 -- Use as overflow to complete your
                answers ]
              \end{center}

\end{questions}

\end{document}
